\subsection{与 Coxeter 系统的关系}

定理 1. 设 \(\mathbf{C}\) 是一个室，令 \(\mathbf{S}\) 为关于 \(\mathbf{C}\) 的墙的反射所成的集合。

\begin{enumerate}
\def\labelenumi{(\roman{enumi})}
\item
有序对 \((\mathbf{W},\mathbf{S})\) 是一个 Coxeter 系统。
\item
设 \(w \in \mathrm{W}\)，且 \(\mathrm{H}\) 是 \(\mathrm{C}\) 的墙。关系 \(l(s_{\mathrm{H}}w) > l(w)\) 蕴含室 \(\mathrm{C}\) 与 \(w(\mathrm{C})\) 位于 \(\mathrm{H}\) 的同一侧。
\item
对于任意室 \(\mathbf{C}'\)，存在唯一的 \(w \in \mathbf{W}\)，使得 \(w(\mathbf{C}) = \mathbf{C}'\)。
\item
满足 \(\mathbf{H}\) 的超平面 \(s_{\mathrm{H}} \in \mathbf{W}\) 的集合等于 \(\mathfrak{H}\)。
\end{enumerate}

S 的每个元素的阶都是 2，而引理 2 表明 S 生成 W。对于 C 的任一墙 H，记 \(P_{H}\) 为满足下列条件的 \(w \in W\) 所成的集合：室 C 与 \(w(\mathrm{C})\)（它们不与 H 相交）位于 H 的同一侧。我们将验证《代数》第 IV 章 §1 第 7 号的 \((\mathrm{A}')\)、\((\mathrm{B}')\) 和 (C) 条件。

\((A') 1 \in P_H\)：显然。

(B') \(P_{H}\) 与 \(s_{H}.P_{H}\) 不相交：事实上，\(w(\mathrm{C})\) 和 \(s_{H}w(\mathrm{C})\) 位于 H 的相反两侧；因此，若 \(w(\mathrm{C})\) 与 C 在 H 的同一侧，它就不会与 \(s_{\mathrm{H}}w(\mathrm{C})\) 在同一侧。

\begin{enumerate}
\def\labelenumi{(\Alph{enumi})}
\setcounter{enumi}{2}
\tightlist
\item
设 \(w \in \mathrm{P_H}\)，且 \(\mathrm{H}'\) 是 \(\mathbf{C}\) 的一面墙，使 \(ws_{\mathrm{H}'} \notin \mathrm{P_H}\)；则 \(ws_{\mathrm{H}'} = s_{\mathrm{H}}w\)：根据假设，\(w(\mathbf{C})\) 与 \(\mathbf{H}\) 在 \(\mathbf{C}\) 的同一侧，而 \(ws_{\mathrm{H}'}(\mathbf{C})\) 在另一侧。因此，\(ws_{\mathrm{H}'}(\mathbf{C})\) 与 \(w(\mathbf{C})\) 位于 H 的相反两侧；所以，\(s_{\mathrm{H}^{\prime}}(\mathrm{C})\) 与 C 位于超平面 \(w^{-1}(\mathrm{H})\) 的相反两侧。令 a 为支撑为 \(H^{\prime}\) 的 C 的一个面上的点。点 \(a = s_{\mathrm{H}^{\prime}}(a)\) 同时属于室 C 与 \(s_{\mathrm{H}^{\prime}}(\mathrm{C})\) 的闭包，而这两个室分别包含于 \(w^{-1}(\mathrm{H})\) 所界定的两个开半空间中；于是 \(a \in w^{-1}(\mathrm{H})\)，故 \(\mathrm{H}^{\prime} = w^{-1}(\mathrm{H})\)。由此推出 \(s_{\mathrm{H}^{\prime}} = w^{-1}s_{\mathrm{H}}w\)，从而 \(ws_{H^{\prime}} = s_{H}w\)。
\end{enumerate}

断言 (i) 和 (ii) 由此以及《代数》第 IV 章 §1 第 7 号命题 6 得出。此外，我们有（同上，条件 (A)）

\[
\bigcap_ {H \in \mathfrak {M}} P _ {H} = \{1 \}.\tag{3}
\]

引理 2 表明 W 在室的集合上传递。此外，若 \(w \in W\) 满足 \(w(\mathrm{C}) = \mathrm{C}\)，则对 C 的每一面墙 H 都有 \(w \in P_{H}\)，因而由 (3) 得 w = 1。这证明了 (iii)。

最后，设 H 是满足 \(s_{H} \in W\) 的一个超平面。如果 H 不属于 H，则至少存在一个与 H 相交的室 \(C'\)（\(\S1\)，第 3 号命题 7）。\(H \cap C'\) 中的每个交点都在 \(s_{H}\) 作用下不变，因而同时属于室 \(C'\) 和 \(s_{\mathrm{H}}(\mathrm{C}')\)；所以，\(\mathrm{C}' = s_{\mathrm{H}}(\mathrm{C}')\)，这与 (iii) 矛盾，因为 \(s_{H} \neq 1\)。

推论. 设 \(\Sigma\) 是生成 W 的一组反射。那么属于 W 的每个反射都与 \(\Sigma\) 中的一个元素共轭。

令 \(\mathfrak{H}'\) 为形如 \(w(\mathrm{H})\) 的超平面所成的集合，其中 \(w \in \mathrm{W}\)、\(\mathrm{H} \in \mathfrak{H}\)，且 \(s_{\mathrm{H}} \in \Sigma\)。由于 \(\mathrm{W}\) 由族 \((s_{\mathrm{H}})_{\mathrm{H} \in \mathfrak{H}'}\) 生成，且 \(\mathfrak{H}'\) 在 \(\mathrm{W}\) 下稳定，我们可以用 \(\mathfrak{H}'\) 代替 \(\mathfrak{H}\) 应用本号的全部结果；但是定理 1(iv) 表明 \(\mathrm{W}\) 中的每个反射都形如 \(s_{\mathrm{H}}\)，其中 \(\mathrm{H} \in \mathfrak{H}'\)，因此推论得证。

